\chapter{设计输入与总体方案}\label{chap:design}

\section{总体设计思路}

系统采用“公共仿真基础、可替换编码模块、场景化信道与统计配置”的组织方式。公共仿真模块（Common）位于三类编码模块之外，提供公共数据结构、BPSK、基础 AWGN、硬判决与 LLR、帧样本池、噪声样本池、信息比特的误码统计、停止控制、断点状态保存、分片合并及公共结果字段。该模块不生成 BCH、卷积码或 LDPC 校验关系，也不执行三类码的核心译码。

BCH、CC 和 LDPC 模块接收相同类型的信息比特和公共控制信息，但各自负责码长适配、编码、译码、信息比特恢复、算法状态以及任务专用指标。多径、频偏、遮挡和突发错误等扩展信道也由具体任务配置。

\section{数据对象与长度关系}

一帧电文经过适配、编码、调整和恢复时会出现多个长度。表~\ref{tab:data-length-objects}给出报告中的统一名称。

\begin{table}[htbp]
  \centering
  \small
  \caption{数据对象、长度符号及附加位关系}
  \label{tab:data-length-objects}
  \begin{tabular}{p{0.20\textwidth}p{0.20\textwidth}p{0.38\textwidth}p{0.14\textwidth}}
    \toprule
    符号 & 数据对象 & 含义 & 附加位关系 \\
    \midrule
    $K_{\mathrm{payload}}$ & 原始有效电文 & 用户需要恢复并用于 BER/FER 的比特 & 不含附加位 \\
    $K_{\mathrm{codec_input}}$ & 编码器输入 & 编码器实际接收的信息侧序列 & 可含 filler、CRC 或适配位 \\
    $N_{\mathrm{encoded}}$ & 编码输出 & Common 码率分母对应的最终编码长度 & 含校验、尾比特等 \\
    $N_{\mathrm{transmitted}}$ & 实际发送序列 & 映射到 BPSK 并进入信道的比特数 & 不含未发送缩短位和打孔位 \\
    $N_{\mathrm{received}}$ & 接收观测 & 进入硬判决或软信息计算的有效观测数 & 受接收处理影响 \\
    $K_{\mathrm{recovered}}$ & 恢复的信息比特 & 去除附加位后提交统计的比特序列 & 应等于 $K_{\mathrm{payload}}$ \\
    \bottomrule
  \end{tabular}
\end{table}

Filler 用于把短的信心比特补到编码结构要求的输入规模，不增加码率；CRC 用于检错或停止判断，也不属于信息比特；卷积码 tail 把编码器状态终止到约定状态，会增加编码输出长度；puncturing 从母码输出中删除部分比特，使实际码率提高；shortening 则把已知的信息位从发送序列中省略，接收端译码前按约定恢复这些位置。相同的 $K_{\mathrm{payload}}$ 因此可能对应不同的 $K_{\mathrm{codec\ input}}$、$N_{\mathrm{encoded}}$ 和 $N_{\mathrm{transmitted}}$。

\section{Common与编码模块的职责分工}

表~\ref{tab:common-codec-responsibility}按数据流列出 Common 和编码任务的边界。公共接口的存在表示三类码可以接入同一处理框架，不表示具体任务已经使用完全相同的信道、统计字段或停止参数。

\begin{table}[htbp]
  \centering
  \scriptsize
  \caption{Common 与三类编码任务的职责分工}
  \label{tab:common-codec-responsibility}
  \begin{tabular}{p{0.15\textwidth}p{0.35\textwidth}p{0.40\textwidth}}
    \toprule
    环节 & Common & 编码任务 \\
    \midrule
    信息比特的输入 & 确定性生成或从 Frame Pool 按帧索引读取 & 选择 200/300 bit 业务长度 \\
    码长适配 & 提供长度结构和公共字段 & 实现 filler、CRC、shortening、tail、puncturing \\
    编码 & 提供编码器接口 & 实现 BCH、卷积码或 BG2 Direct LDPC \\
    调制 & 固定实值 BPSK，0 映射 $+1$，1 映射 $-1$ & 通常直接复用 \\
    信道 & AWGN 实现和信道接口 & 增加多径、频偏、遮挡、突发错误等 \\
    解调 & 硬判决和未量化 LLR & 设置量化、裁剪、饱和或补偿 \\
    译码 & 提供输入/输出接口 & 实现查表、代数、Viterbi、BP 或 NMS \\
    信息比特的恢复 & 提供统计入口和长度校验 & 去除 filler、CRC、tail 等附加位 \\
    统计 & BER/FER、基础成功率、平均和最大时间 & 增加误纠、声明成功、迭代、复杂度、分位数等 \\
    \bottomrule
  \end{tabular}
\end{table}

\section{低速BCH处理路径}

低速路径从 \SI{200}{bit} 或 \SI{300}{bit} 信息比特开始。分组方案先按 BCH 信息位长度切分，不足一组时补入填充；整块方案把完整信息比特映射到较长 BCH 母码并执行缩短。编码产生一个或多个系统 BCH 码字，S7 可在编码后加入交织，随后进行 BPSK 和信道传输。接收端先硬判决，再在启用交织时去交织，之后执行查表或代数 BCH 译码，删除补入或缩短对应的非信息比特位置，最后提交固定长度的恢复结果。

分组方案会产生多个短码字，一帧 FER 的判定取决于所有子块恢复后的信息比特：只要任一信息比特位错误，整帧计为错误帧。整块方案的纠错边界、缩短位置和误纠状态由 BCH 章节结合实际实现说明。

\section{高速卷积码处理路径}

高速卷积码输入固定为 \SI{300}{bit}信息比特。整块组织可在电文末尾加入零尾，使移位寄存器回到约定状态；按时隙处理可保留连续状态，并通过滑窗 Viterbi 降低等待完整电文的时间。项目以 $1/2$ 母码为基础，通过打孔形成接近 $2/3$、$3/4$ 的发送码率。

接收端可将符号硬判为 0/1，也可把实值观测转换为 LLR 或量化值。硬、软 Viterbi 共用相同码字和接收符号时，差异主要来自分支度量；回溯深度、滑窗长度、软值位宽和饱和策略会同时影响存储、首输出和译码时间。整块零尾、连续状态和时隙分块还会改变结构时延。

\section{高速LDPC处理路径}

LDPC 路径只处理 \SI{300}{bit} 整块信息比特。编码模块依据 BG2 短块结构选择提升因子和使用的子矩阵，通过 filler 或 shortening 适配编码输入，并由 Direct QC-LDPC 编码器生成校验位。接收端计算软信息后，分别使用 Layered BP 和 Layered NMS 迭代译码；两者设置最大迭代次数，并在校验方程全部满足时允许提前停止。恢复阶段删除 filler 等非信息比特位置，只把原始 \SI{300}{bit} 范围交给公共统计。

本研究采用基于 5G NR BG2 结构的短块实现，其适用范围不包含完整的 5G NR 物理层传输链。

\section{统一通信与仿真链路}

三类码共享的上层数据流为：信息比特、码长适配、编码、可选交织、BPSK调制、信道、硬判决、去交织、译码、信息比特恢复和统计，如图\ref{fig:overall-scenario}所示。
% VISIO-ABSOLUTE-PATH: C:/Users/V3169/Desktop/Project/SCL_Channel_Coding_Design/Report/figures/visio/common/overall_scenario_architecture.png
\begin{figure}[htbp]
  \centering
  \includegraphics[width=0.90\textwidth]
  {figures/visio/common/overall_scenario_architecture.pdf}
  \caption{总体编码与仿真架构图}
  \label{fig:overall-scenario}
\end{figure}

结果整理保留内部 \texttt{ebN0Db} 和实际码率，并按每个点的码率计算正式横轴符号信噪比。该处理只增加展示坐标，不改动接收符号、错误计数或停止结果。下层控制信息通过帧索引、样本池身份和配置标识与各仿真点关联。

\begin{table}[htbp]
  \centering
  \small
  \caption{三类编码处理路径的主要组织差异}
  \label{tab:codec-path-comparison}
  \begin{tabular}{p{0.17\textwidth}p{0.23\textwidth}p{0.24\textwidth}p{0.26\textwidth}}
    \toprule
    项目 & BCH & 卷积码 & BG2 Direct LDPC \\
    \midrule
    业务输入 & 200/300 bit 低速 & 300 bit 高速 & 300 bit 高速 \\
    组织 & 分组或整块缩短 & 整块零尾或时隙连续 & 整块短块适配 \\
    码率调整 & 码型选择、缩短、填充 & 母码、尾比特、打孔 & 提升因子、filler、Direct 结构 \\
    译码输入 & 硬判决 & 硬比特或量化软信息 & LLR 或量化软信息 \\
    译码状态 & 无错、纠错、失败、误纠 & 路径输出及任务状态 & 校验满足、迭代上限等 \\
    交织 & S7 可选 & S7 可选 & 禁止配置 \\
    关键时延 & 子块处理与整帧完成 & 首输出、回溯、滑窗 & 迭代次数和尾部时间 \\
    \bottomrule
  \end{tabular}
\end{table}

\section{本章小结}

总体方案将公共数据与控制链同三类编译码实现分离：公共仿真模块固定可复用的基础调制、AWGN、随机样本、统计和一致性机制，编码模块负责码长适配、编译码及专用状态。低速 BCH、高速卷积码和高速LDPC因组织方式不同保留独立长度和时延字段；S1--S7依次完成码内筛选、场景检验、实现代价比较和交织分析。
